\subsection{普遍可测函数}

在本节中，X 是局部紧拓扑空间。我们回顾（INT, V, p. 28, § 3, n° 4，定义 2），从 X 到拓扑空间 Y 中的映射 f 称为普遍可测的，如果它对 X 上的每个正测度 \(\mu\) 都是 \(\mu\)-可测的。为此只需 f 对 X 上每个具有紧支撑的正测度 \(\mu\) 是 \(\mu\)-可测的（INT, V, p. 28, § 4, no. 3，命题 6）。

引理 4. --- 设 Y 与 Z 是局部紧拓扑空间，\(f\colon X\to Y\) 与 \(g\colon Y\to Z\) 是普遍可测映射。映射 \(g\circ f\colon X\to Z\) 是普遍可测的。

设 \(\mu\) 是 X 上具有紧支撑的正测度，C 是它的支撑。\(f\) 在 C 上的限制是 \((\mu|\mathrm{C})\)-正常的。映射 \(g\) 关于像测度 \((f|\mathrm{C})(\mu|\mathrm{C})\) 可测，故映射 \((g \circ f)|\mathrm{C} = g \circ (f|\mathrm{C})\) 是 \((\mu|\mathrm{C})\)-可测的。因此映射 \(g \circ f\) 是 \(\mu\)-可测的。引理得证。

我们用 \(\mathcal{L}_{\mathrm{u}}(\mathrm{X})\) 表示在 X 上普遍可测的复值函数所成的向量空间，用 \(\mathcal{L}_{\mathrm{u}}^{\infty}(\mathrm{X})\) 表示在 X 上有界的函数 \(f \in \mathcal{L}_{\mathrm{u}}(\mathrm{X})\) 所成的子空间。对 \(f \in \mathcal{L}_{\mathrm{u}}^{\infty}(\mathrm{X})\)，我们记 \(\|f\|_{\infty} = \sup_{x \in X} |f(x)|\) 。

命题 3. --- a) 集 \(\mathcal{L}_{\mathrm{u}}(\mathrm{X})\) 是从 X 到 C 中的函数所成的对合代数的对合子代数；

\begin{enumerate}
\def\labelenumi{\alph{enumi})}
\setcounter{enumi}{1}
\item
  集 \(\mathcal{L}_{\mathrm{u}}^{\infty}(\mathrm{X})\) 是 \(\mathcal{L}_{\mathrm{u}}(\mathrm{X})\) 的子代数，且由 \(f\mapsto \| f\|_{\infty}\) 定义的映射是 \(\mathcal{L}_{\mathrm{u}}^{\infty}(\mathrm{X})\) 上的一个范数，对此范数 \(\mathcal{L}_{\mathrm{u}}^{\infty}(\mathrm{X})\) 是有单位元的星代数；
\item
  代数 \(\mathcal{L}_{\mathrm{u}}(\mathrm{X})\) 包含 \(\mathcal{C}(\mathrm{X})\) ，且 \(\mathcal{L}_{\mathrm{u}}^{\infty}(\mathrm{X})\) 包含 \(\mathcal{C}_b(\mathrm{X})\) ；
\item
  从 X 到 \(\mathbf{C}\) 中的每个博雷尔函数（TG, IX, p. 61，定义 5）都属于 \(\mathcal{L}_{\mathrm{u}}(\mathrm{X})\) ；
\item
  对任意序列 \((f_n)_{n\in \mathbf{N}}\) ——\(\mathcal{L}_{\mathrm{u}}(\mathrm{X})\) 中逐点收敛到函数 \(f\) （从 X 到 \(\mathbf{C}\) 中的函数）的序列——我们有 \(f\in \mathcal{L}_{\mathrm{u}}(\mathrm{X})\) 。
\end{enumerate}

断言 \(a\)）与 \(c\)）由定义得出（参见 INT, IV, p. 175, § 5, no. 3，推论 3）。断言 \(d\)）是 INT, IV, p. 179, § 5, no. 5 定理 4 的推论，因为 \(\mathbf{C}\) 的每个博雷尔子集在博雷尔函数下的原像是 X 的博雷尔子集，从而是普遍可测的（INT, V, p. 28, § 3, no. 4）。最后，断言 \(e\)）由叶戈罗夫定理得出（INT, IV, p. 175, § 5, no. 4，定理 2）。

为证明断言 \(b\)），只需注意断言 \(e\)）蕴含，\(a\) 更不用说，普遍可测函数的一致极限是普遍可测的。

X 上可能存在不是博雷尔的普遍可测函数。事实上，若 X 可度量化，X 的苏斯林子集的特征函数是普遍可测的（参见 INT, IV, p. 171, § 5, no. 1，推论 2）；然而在每个不可数的波兰空间中，都存在不是博雷尔的苏斯林子集（“苏斯林定理”；这由 TG, IX, p. 120 习题 8 以及对每个不可数波兰空间 X 存在逆映射为博雷尔的双射 \(\mathrm{X}\to\mathbf{N}^{\mathbf{N}}\) 这一事实得出）。

引理 5. --- 设 \(\mu\) 是 X 上的正测度。设 \(g\) 是 \(\mu\)-可测函数（从 X 到 \(\mathbf{C}\) 中的），并设 \(S\) 是它的 \(\mu\)-本质像。对每个函数 \(f \in \mathcal{L}_{\mathrm{u}}(\mathrm{S})\)，函数 \(h\) （从 X 到 \(\mathbf{C}\) 中的函数），若满足 \(h(x) = 0\) （当 \(g(x) \notin \mathrm{S}\)）与 \(h(x) = f(g(x))\) （当 \(g(x) \in \mathrm{S}\)），则是 \(\mu\)-可测的。

设 \(x \in X\) ，\(U\) 是 x 的相对紧开邻域。集 \(\mathrm{N} = \mathrm{U} - (\mathrm{U} \cap g^{-1}(\mathrm{S}))\) 是 \((\mu|\mathrm{U})\)-零测的。对每个整数 \(n \geqslant 1\)，设 \(V_{n}\) 是满足 \(N \subset V_{n} \subset U\) 且 \(\mu(V_{n} - N) < 1/n\) 的开集（INT, IV, p. 116, § 1, no. 4，命题 19）。

限制 \(\widetilde{g}\) （\(g\) 在 \(\mathrm{U}-\mathrm{V}_n\) 上的限制）是 \(\mu|(\mathrm{U}-\mathrm{V}_n)\)-正常的，且其像含于 S。由于 \(f\) 普遍可测，映射 \(x\mapsto f(g(x))\) （从 \(\mathrm{U}-\mathrm{V}_n\) 到 \(\mathbf{C}\) 中的映射）关于测度 \(\mu|(\mathrm{U}-\mathrm{V}_n)\) 可测（INT, V, p. 71, § 6, no. 2，命题 3）。因此，映射 \(h_n\) （从 X 到 \(\mathbf{C}\) 中的映射）——若 \(h_n(x) = 0\) （当 \(x\notin \mathrm{U}-\mathrm{V}_n\)）且 \(h_n(x) = f(g(x))\) （当 \(x\in \mathrm{U}-\mathrm{V}_n\)）——是 \(\mu\)-可测的（INT, IV, p. 193, § 5, no. 10，命题 16）。

由于 \(h_{n}(x) \to h(x)\) 对 \(\mu\)-几乎所有的 \(x \in U\) 成立，\(h\) 在 U 上的限制是 \(\mu\)-可测的（INT, IV, p. 175, § 5, no. 4，定理 2）。于是可断定 \(h\) 是 \(\mu\)-可测的，依据是局部化原理（INT, IV, p. 171, § 4, no. 2，命题 4）。

注. --- 在该引理的假设下，我们约定滥用记号，用 \(f \circ g\) 表示引理中定义的函数 \(h\) 。

若 \(g = g_{1}\) 在 X 上局部 \(\mu\)-几乎处处成立，则函数 \(g\) 与 \(g_{1}\) 有相同的 \(\mu\)-本质像，且有 \(f \circ g = f \circ g_{1}\) 。

若 \(g\) 是在 X 上 \(\mu\)-几乎处处定义的 \(\mu\)-可测函数，我们也用 \(f \circ g\) 表示函数 \(f \circ g_{1}\) ，其中 \(g_{1}\) 是定义在 X 上且局部 \(\mu\)-几乎处处等于 g 的 \(\mu\)-可测函数。

